%%%%%%%% OSDI 2026 PAPER TEMPLATE %%%%%%%%%%%%%%%%%
%
% The 20th USENIX Symposium on Operating Systems Design and Implementation
% July 13--15, 2026, Seattle, WA, USA
%
% Format: <= 12 pages (excluding references), 8.5"x11", 10pt on 12pt leading,
%         two-column, Times Roman, 7"x9" text block
% Camera-ready: <= 14 pages (2 extra pages allowed)
%
% Official CFP: https://www.usenix.org/conference/osdi26/call-for-papers
% Template source: https://www.usenix.org/conferences/author-resources/paper-templates
%
% IMPORTANT NOTES:
% - OSDI 2026 has two tracks: Research and Operational Systems
% - For Operational Systems track, title must end with "(Operational Systems)"
% - Max 8 submissions per author
% - Papers should be the right length (not padded to 12 pages)
% - Papers <= 6 pages are unlikely to receive full consideration
% - Use anonymized project/system name (different from arXiv/talks)
%
% WHAT OSDI REVIEWERS LOOK FOR:
% 1. Significant problem motivation
% 2. Interesting and compelling solution
% 3. Practicality and benefits demonstrated
% 4. Clear contribution articulation
% 5. Advances beyond previous work

\documentclass[letterpaper,twocolumn,10pt]{article}
\usepackage{usenix-2020-09}

% Recommended packages for systems papers
\usepackage[utf8]{inputenc}
\usepackage{amsmath,amssymb}
\usepackage{graphicx}
\usepackage{booktabs}       % Professional tables
\usepackage{hyperref}
\usepackage{url}
\usepackage{xspace}
\usepackage{subcaption}     % Side-by-side figures
\usepackage{algorithm}      % Algorithm environment
\usepackage{algorithmic}    % Pseudocode formatting
\usepackage{listings}       % Code listings
\usepackage[capitalize,noabbrev]{cleveref}  % Smart cross-references

% Code listing style for systems papers
\lstset{
  basicstyle=\footnotesize\ttfamily,
  numbers=left,
  numberstyle=\tiny,
  xleftmargin=2em,
  breaklines=true,
  tabsize=2,
  showstringspaces=false,
  frame=single,
  captionpos=b
}

% Custom commands -- replace \system with your anonymized name
\newcommand{\system}{SystemName\xspace}
\newcommand{\eg}{e.g.,\xspace}
\newcommand{\ie}{i.e.,\xspace}
\newcommand{\etal}{\textit{et al.}\xspace}
\newcommand{\para}[1]{\smallskip\noindent\textbf{#1.}}
\newcommand{\parait}[1]{\smallskip\noindent\textit{#1.}}

% Convenience macros for units (common in systems papers)
\newcommand{\us}{\,$\mu$s\xspace}
\newcommand{\ms}{\,ms\xspace}
\newcommand{\GB}{\,GB\xspace}
\newcommand{\MB}{\,MB\xspace}
\newcommand{\Gbps}{\,Gbps\xspace}

\begin{document}

% For submission: use anonymized title and Paper #XXX as author
% For Operational Systems track: add "(Operational Systems)" to title
\title{Your Paper Title Here}
% \title{Your Paper Title Here (Operational Systems)}  % Operational Systems track

\author{Paper \#XXX}  % Anonymized for submission
% Camera-ready:
% \author{
%   {\rm Author One}\\
%   Affiliation One\\
%   \texttt{email@example.com}
%   \and
%   {\rm Author Two}\\
%   Affiliation Two\\
%   \texttt{email@example.com}
% }

\maketitle

%----------------------------------------------------------------------
\begin{abstract}
% Guidelines for a strong OSDI abstract:
% - State what you achieved (the contribution)
% - Why this is hard and important (the problem)
% - How you do it (the approach)
% - What evidence you have (evaluation highlights)
% - Your most remarkable result (the hook)
%
% Keep to 150--200 words. Avoid citations in the abstract.

We present \system, a [describe system] that [key capability].
[Problem statement: why existing approaches fall short.]
\system addresses this through [key technique/insight].
We evaluate \system on [workloads/benchmarks] and show that it achieves
[X]\% improvement in [metric] over [baseline], while reducing [other metric]
by [Y]$\times$.
\end{abstract}

%----------------------------------------------------------------------
\section{Introduction}
\label{sec:intro}

% Structure your introduction as follows:
% 1. Problem context and motivation (1--2 paragraphs)
% 2. Why existing solutions are insufficient (1 paragraph)
% 3. Key insight / approach overview (1 paragraph)
% 4. Contributions (bulleted list)
% 5. Results highlights (1 paragraph)
%
% OSDI reviewers look for: significant problem + compelling solution +
% demonstrated practicality + clear contributions + advances beyond prior work.

Modern systems face the challenge of [describe problem]~\cite{dean2004mapreduce}.
As workloads grow in scale and complexity, existing approaches such as
[prior work]~\cite{abadi2016tensorflow} struggle to [limitation].

\para{Key Insight}
Our key observation is that [insight]. This enables \system to [capability]
without [drawback of prior approaches].

We make the following contributions:
\begin{itemize}
  \item We identify [problem/opportunity] and characterize its impact
        on [workloads] (\cref{sec:background}).
  \item We design \system, which introduces [technique] to address
        [challenge] (\cref{sec:design}).
  \item We implement \system in [X] lines of [language] and integrate
        it with [existing system] (\cref{sec:implementation}).
  \item We evaluate \system on [benchmarks] and demonstrate [X]$\times$
        improvement over [baseline] (\cref{sec:evaluation}).
\end{itemize}

%----------------------------------------------------------------------
\section{Background and Motivation}
\label{sec:background}

% Provide context the reader needs to understand your contribution.
% Include a motivating example or measurement study.

\subsection{Problem Context}

Describe the system context and relevant background.
Prior work~\cite{moritz2018ray, zaharia2012spark} has explored [related area],
but [gap remains].

\subsection{Motivating Example}

% Use concrete numbers from real workloads to motivate the problem.
\Cref{fig:motivation} shows [measurement] across [workloads].
We observe that [finding], which motivates our approach.

\begin{figure}[t]
  \centering
  % Replace with your actual figure
  \fbox{\parbox{0.9\columnwidth}{\centering\vspace{3em}
    \textit{Motivating measurement or characterization study}
  \vspace{3em}}}
  \caption{[Description of motivating measurement.] We observe that
    [key finding] across [N] production workloads, motivating the need
    for [your approach].}
  \label{fig:motivation}
\end{figure}

%----------------------------------------------------------------------
\section{Design}
\label{sec:design}

% Present your system design top-down:
% 1. Architecture overview (with figure)
% 2. Key components / mechanisms
% 3. How they interact

\Cref{fig:architecture} shows the overall architecture of \system.
The system consists of [N] key components: [list].

\begin{figure*}[t]
  \centering
  % Replace with your actual architecture diagram
  \fbox{\parbox{0.9\textwidth}{\centering\vspace{4em}
    \textit{System architecture diagram showing key components and data flow}
  \vspace{4em}}}
  \caption{Architecture of \system. [Component A] handles [function],
    while [Component B] manages [function]. Arrows indicate [data/control flow].}
  \label{fig:architecture}
\end{figure*}

\subsection{Component A: [Name]}

Describe the first key component. A formal specification of the core
algorithm can be found in \cref{alg:core}.

\begin{algorithm}[t]
  \caption{Core algorithm of \system}
  \label{alg:core}
  \begin{algorithmic}[1]
    \STATE \textbf{Input:} workload $W$, resources $R$
    \STATE \textbf{Output:} scheduling plan $P$
    \STATE Initialize plan $P \leftarrow \emptyset$
    \FOR{each task $t_i \in W$}
      \STATE Estimate resource demand $d_i \leftarrow \text{Predict}(t_i)$
      \IF{$\text{Available}(R) \geq d_i$}
        \STATE $P \leftarrow P \cup \{(t_i, \text{Allocate}(R, d_i))\}$
      \ELSE
        \STATE Enqueue $t_i$ for deferred scheduling
      \ENDIF
    \ENDFOR
    \STATE \textbf{return} $P$
  \end{algorithmic}
\end{algorithm}

\subsection{Component B: [Name]}

Describe the second key component. The expected throughput can be
modeled as:
\begin{equation}
  \label{eq:throughput}
  T = \frac{N \cdot B}{L + \frac{B}{C}}
\end{equation}
where $N$ is the number of parallel workers, $B$ is the batch size,
$L$ is the network latency, and $C$ is the per-worker compute rate.

\subsection{Handling Edge Cases}

Discuss how the design handles failures, stragglers, or other
edge cases important in production systems.

%----------------------------------------------------------------------
\section{Implementation}
\label{sec:implementation}

% Describe implementation details that matter for reproducibility.
% Include system size, language, key libraries, and integration points.

We implement \system in approximately [X]K lines of [language].
Key implementation details include:

\para{Threading Model}
[Describe the threading/concurrency model.]

\para{Integration}
\system integrates with [existing system] by [method of integration].
We modify [N] lines of the original codebase.

%----------------------------------------------------------------------
\section{Evaluation}
\label{sec:evaluation}

% Structure your evaluation to answer specific questions:
% - Q1: How does \system compare to state-of-the-art? (end-to-end)
% - Q2: What is the contribution of each component? (ablation)
% - Q3: How does \system scale? (scalability)
% - Q4: What is the overhead? (cost analysis)

We evaluate \system to answer the following questions:
\begin{enumerate}
  \item How does \system compare to state-of-the-art systems?
  \item What is the contribution of each design component?
  \item How does \system scale with increasing workload?
  \item What overhead does \system introduce?
\end{enumerate}

\subsection{Experimental Setup}
\label{sec:eval:setup}

\para{Testbed}
We run experiments on a cluster of [N] machines, each with
[CPU model], [X]\GB RAM, and [GPU model if applicable],
connected via [network].

\para{Workloads}
We use [N] workloads from [source]: [list workloads].
\Cref{tab:workloads} summarizes their characteristics.

\para{Baselines}
We compare against [N] baselines:
(1)~[Baseline A]~\cite{verma2015borg},
(2)~[Baseline B]~\cite{ongaro2014raft}, and
(3)~[Baseline C].

\begin{table}[t]
  \caption{Workload characteristics used in evaluation.
    [Describe what the columns represent.]}
  \label{tab:workloads}
  \centering
  \begin{small}
  \begin{tabular}{@{}lrrrl@{}}
    \toprule
    \textbf{Workload} & \textbf{Tasks} & \textbf{Data (GB)} &
    \textbf{Duration} & \textbf{Type} \\
    \midrule
    WorkloadA  & 1,024  & 128  & 2.4\,h  & Batch     \\
    WorkloadB  & 512    & 64   & 1.1\,h  & Streaming \\
    WorkloadC  & 4,096  & 512  & 8.7\,h  & ML Train  \\
    WorkloadD  & 256    & 32   & 0.5\,h  & Serving   \\
    \bottomrule
  \end{tabular}
  \end{small}
\end{table}

\subsection{End-to-End Performance}
\label{sec:eval:e2e}

\Cref{tab:e2e} shows the end-to-end performance comparison.
\system achieves [X]\% higher throughput and [Y]\% lower latency
compared to [best baseline].

\begin{table}[t]
  \caption{End-to-end performance comparison. Bold indicates best result.
    \system achieves the highest throughput and lowest p99 latency
    across all workloads.}
  \label{tab:e2e}
  \centering
  \begin{small}
  \begin{tabular}{@{}lccc@{}}
    \toprule
    \textbf{System} & \textbf{Throughput} & \textbf{p50 Latency} &
    \textbf{p99 Latency} \\
                     & \textbf{(Kops/s)}  & \textbf{(ms)}        &
    \textbf{(ms)}        \\
    \midrule
    Baseline A       & 125.3               & 4.2                  & 18.7 \\
    Baseline B       & 142.1               & 3.8                  & 15.2 \\
    Baseline C       & 98.6                & 5.1                  & 22.4 \\
    \textbf{\system} & \textbf{187.4}      & \textbf{2.9}         & \textbf{9.8} \\
    \bottomrule
  \end{tabular}
  \end{small}
\end{table}

\subsection{Ablation Study}
\label{sec:eval:ablation}

To understand the contribution of each component, we evaluate variants
of \system with individual components disabled.
\Cref{fig:ablation} shows the results.

\begin{figure}[t]
  \centering
  % Replace with your ablation study figure
  \fbox{\parbox{0.9\columnwidth}{\centering\vspace{3em}
    \textit{Bar chart: \system vs.\ variants with components disabled}
  \vspace{3em}}}
  \caption{Ablation study results. Each bar represents \system with
    one component removed. Component A contributes [X]\% and
    Component B contributes [Y]\% of the total improvement.}
  \label{fig:ablation}
\end{figure}

\subsection{Scalability}
\label{sec:eval:scale}

We evaluate how \system scales from [N] to [M] nodes.
As shown in \cref{fig:scalability}, \system achieves near-linear
scaling up to [K] nodes.

\begin{figure}[t]
  \centering
  \fbox{\parbox{0.9\columnwidth}{\centering\vspace{3em}
    \textit{Line chart: throughput vs.\ number of nodes for each system}
  \vspace{3em}}}
  \caption{Scalability comparison. \system achieves [X]\% of ideal
    linear scaling at [K] nodes, compared to [Y]\% for [baseline].}
  \label{fig:scalability}
\end{figure}

%----------------------------------------------------------------------
\section{Discussion}
\label{sec:discussion}

% Discuss limitations, lessons learned, and generalizability.
% OSDI reviewers appreciate honest discussion of limitations.

\para{Limitations}
[Discuss known limitations of your system.]

\para{Lessons Learned}
[Share insights from building and deploying the system.]

%----------------------------------------------------------------------
\section{Related Work}
\label{sec:related}

% Organize by theme, NOT paper-by-paper.
% Clearly distinguish your work from each category.

\para{[Category A] Systems}
Prior work on [category]~\cite{dean2004mapreduce, abadi2016tensorflow}
focuses on [aspect]. \system differs by [distinction].

\para{[Category B] Approaches}
[Other approaches]~\cite{lamport1978time, verma2015borg} address [problem]
through [method]. In contrast, \system [distinction].

%----------------------------------------------------------------------
\section{Conclusion}
\label{sec:conclusion}

We presented \system, a [description] that [key capability].
Through [technique], \system achieves [improvement] over
state-of-the-art systems. Our evaluation on [workloads] demonstrates
[key results]. [Optional: future work direction.]

%----------------------------------------------------------------------
% Bibliography
% USENIX uses the acm bibliography style
{\footnotesize \bibliographystyle{acm}
\bibliography{references}}

%----------------------------------------------------------------------
% Optional: Appendix (after bibliography for USENIX)
% \appendix
% \section{Additional Evaluation Results}
% Include supplementary material here.

\end{document}
